\documentclass[11pt]{amsart}
\usepackage[margin=1in]{geometry}
\usepackage{amsmath,amssymb,amsthm,mathtools}
\usepackage{booktabs}
\usepackage{enumitem}
\usepackage[colorlinks=true,linkcolor=blue,citecolor=blue,urlcolor=blue]{hyperref}

\newtheorem{theorem}{Theorem}[section]
\newtheorem{lemma}[theorem]{Lemma}
\newtheorem{proposition}[theorem]{Proposition}
\newtheorem{corollary}[theorem]{Corollary}
\newtheorem*{mainthm}{Main Theorem}
\newtheorem*{maincor}{Corollary}
\theoremstyle{definition}
\newtheorem{definition}[theorem]{Definition}
\newtheorem{remark}[theorem]{Remark}

\newcommand{\Z}{\mathbb{Z}}
\newcommand{\Q}{\mathbb{Q}}
\newcommand{\C}{\mathbb{C}}
\newcommand{\HH}{\mathbb{H}}
\newcommand{\T}{\mathcal{T}}
\newcommand{\Aut}{\operatorname{Aut}}
\newcommand{\Out}{\operatorname{Out}}
\newcommand{\Inn}{\operatorname{Inn}}
\newcommand{\End}{\operatorname{End}}
\newcommand{\Sym}{\operatorname{Sym}}
\newcommand{\coker}{\operatorname{coker}}
\newcommand{\Gr}{\operatorname{Gr}}
\newcommand{\Rinf}{R_\infty}

\title{Twisted conjugacy in irreducible spherical Artin groups}
\author{[Author names]}
\date{Draft, November 2026}

\begin{document}

\begin{abstract}
We prove that the Artin groups of types $D_5$, $E_6$, $E_7$, $E_8$, $F_4$, $H_3$ and $H_4$ have property $\Rinf$.
Together with previously known cases, this shows that every irreducible Artin group of spherical type and rank $\ge2$ has
$\Rinf$.

The proof uses the characteristic pure subgroup and its abelianization, viewed as the permutation lattice on the
reflections. Rational, absolutely irreducible constituents of multiplicity one yield a fixed vector after composing
with a suitable inner automorphism. The finite action that must be controlled is governed by automorphisms of
Coxeter groups, so no description of $\Aut(A)$ is required. The constituents are explicit: eigenspaces of strongly
regular graphs for $E_n$, frame modules for $F_4$ and $H_3$, and matrix coefficients in the icosian model for $H_4$.

The same argument gives uniform proofs in types $A_n$, $B_n$ and $D_n$; for $D_5$ this case appears to be new.
\end{abstract}

\maketitle

\noindent\fbox{\parbox{\dimexpr\textwidth-2\fboxsep-2\fboxrule}{\small
\textbf{[Contribution statement --- placeholder, to be finalised by the authors.]}
This work was developed with substantial assistance from AI systems (Anthropic's Claude) for the mathematical
argument, the programs, the Lean formalisation and the text. So far it has been reviewed only by AI systems: an
internal audit, two rounds of an AI ``cold referee'', and an external AI review of the manuscript. None of this is journal peer
review, and no human expert has yet checked the work. The authors must replace this paragraph with an accurate
statement of contributions before submission or posting.}}

\section{Introduction}

Let $G$ be a group and $\varphi\in\Aut(G)$. Elements $x,y\in G$ are \emph{$\varphi$-twisted conjugate} if
$y=zx\varphi(z)^{-1}$ for some $z\in G$. The number $R(\varphi)\in\{1,2,\dots,\infty\}$ of classes is the
\emph{Reidemeister number}, and $G$ has \emph{property $\Rinf$} if $R(\varphi)=\infty$ for every
$\varphi\in\Aut(G)$ \cite{FT15,FN16}.

For an Artin group $A=A(\Gamma)$ with Coxeter group $W=W(\Gamma)$, the question ``which Artin groups have $\Rinf$?''
is Problem~8.6 of the AIM problem list on Artin groups \cite{AIMPL}. Known cases include:
\begin{itemize}
\item braid groups \cite{FG10} and pure braid groups \cite{DGO21};
\item the spherical types $A_n$, $B_n$, $D_4$, $I_2(m)$ \cite{CS22};
\item $D_n$ for $n\ge6$ \cite[Thm.~1]{SV24};
\item many non-spherical types \cite{Juh11,SV24,Sor26}.
\end{itemize}
The case $D_5$ was left open in \cite{SV24}. The authors write there: ``The only remaining case in the family
$A[D_n]$ is $A[D_5]$. At present we do not have a description of $\Aut(A[D_5])$ \dots; we conjecture though that
both $A[D_5]$ and $A[D_5]/Z(A[D_5])$ have property $\Rinf$'' \cite[\S1]{SV24}. Paris--Soroko later computed
$\Aut(A[D_5])$ \cite{PS26}, but they do not discuss $\Rinf$. (The survey list in \cite[\S1]{Sor26} attributes
$D_n$, $n\ge5$, to \cite{SV24}; as far as we can see, $D_5$ is not proved there.)
For spherical types the known proofs use an explicit description of $\Aut(A)$ or $\Aut(A/Z(A))$. For $I_2(m)$
they use instead the virtual freeness of $A/Z(A)$ \cite{CS22}.

\begin{mainthm}
The Artin groups of types $D_5$, $E_6$, $E_7$, $E_8$, $F_4$, $H_3$ and $H_4$ have property $\Rinf$.
\end{mainthm}

For $E_6,\dots,H_4$ see Section~\ref{sec:exceptional}. For $D_5$ see Proposition~\ref{prop:BD} and
Remark~\ref{rem:D5}, where we also note a short alternative derivation from \cite[Thm.~1]{PS26}.

\begin{maincor}
Every irreducible Artin group of spherical type and rank $\ge2$ has property $\Rinf$.
\end{maincor}

The corollary follows from the Main Theorem together with \cite{CS22} (types $A_n,B_n,D_4,I_2(m)$) and
\cite[Thm.~1]{SV24} ($D_n$, $n\ge6$). Rank one must be excluded:
$A(A_1)\cong\Z$, and inversion on $\Z$ has Reidemeister number $2$. The method also gives uniform proofs for
$A_n$, $B_n$ and $D_n$ (Section~\ref{sec:classical}), so among the irreducible spherical types only $I_2(m)$ needs
a different argument (Remark~\ref{rem:dihedral}).

\smallskip\noindent\emph{Significance and novelty.} For the six exceptional types $\Aut(A)$ is not known (cf.\
\cite[\S1]{PS26}). The general criteria based on hyperbolic actions of $\Aut(A)$, or on $\Aut$-invariant
quasimorphisms \cite{FFW23,FFIMSW25,SV26}, also require control of $\Aut(A)$.

The Main Theorem would follow from Conjecture~(NA) of \cite[\S6]{SV24}: finitely generated, residually finite,
non-amenable groups have $\Rinf$. This appears as a theorem in \cite{FT12}, but its proof is reported to be
flawed \cite[\S6]{SV24}. Our results verify (NA), and Conjecture~(R) \cite[Problem~19.28]{Kourovka}, for these
groups.

The broad strategy of obtaining an eigenvalue one via representation theory has precedents:
\begin{itemize}
\item Dekimpe--Gon\c{c}alves--Ocampo \cite{DGO21}, for pure braid groups;
\item Dekimpe--De Rock--Penninckx \cite[Thm.~4.7]{DDP09}, for flat manifolds. They require a $\Q$-irreducible
$\Z$-subrepresentation $\eta_0$ of the holonomy representation whose image is not $\Q$-conjugate to that of any
other $\Q$-subrepresentation (in particular, it has multiplicity one), and such that for every $D_0$ in the
integral normaliser of $\eta_0(F)$ some $\eta_0(A)D_0$, $A\in F$, has eigenvalue one;
\item the determinant criterion of Dekimpe--Penninckx for crystallographic groups \cite[Cor.~3.10]{DP11}, as
stated in \cite[Thm.~3.3]{DKT19};
\item the eigenvalue-one property of finite groups announced by Hiss--Lutowski \cite{HL24}. For odd-dimensional
constituents this gives an alternative route to the witnesses.
\end{itemize}
What is new here is:
\begin{itemize}
\item the application to Artin groups through the \emph{characteristic reflection lattice}
$H_1(P;\Z)\cong\Z^\T$;
\item the control of the relevant finite symmetries by $\Aut(W)$ \cite{FH03};
\item the explicit sign and swap witnesses;
\item the completion of the exceptional cases, and of $D_5$.
\end{itemize}

\smallskip\noindent\emph{Overview.} Let $P=\ker(A\to W)$, let $\T$ be the set of reflections, and let
$\rho\colon W\to\Sym(\T)$ be the conjugation action.
\begin{itemize}
\item $P$ is characteristic, and $H_1(P;\Z)\cong\Z^\T$ is a permutation module (Section~\ref{sec:lattice}).
\item An automorphism $\varphi$ acts on $\Z^\T$ as $\pi N$, where $\pi$ realises the induced automorphism of
$\rho(W)\cong W/Z(W)$ and $N$ commutes with $\rho(W)$. On a good constituent $U$, $N=\pm1$.
\item For a suitable inner automorphism $\iota_a$, the composite $\varphi'=\iota_a\circ\varphi$ has a fixed vector.
So $R(\varphi'|_P)=\infty$, hence $R(\varphi')=\infty$, and $R(\varphi)=R(\varphi')$ (Section~\ref{sec:criterion}).
\item The possible $\pi$ are controlled by $\Aut(W)$. Explicit constituents and witnesses finish the proof
(Section~\ref{sec:exceptional}).
\end{itemize}

\smallskip\noindent\emph{Formal verification.} A Lean~4/Mathlib formalisation accompanies this paper. It contains
no \texttt{sorry}, no \texttt{axiom} and no \texttt{native\_decide}. The Lean kernel checks:
\begin{itemize}
\item the Reidemeister lemmas and the criterion (both cases);
\item the Artin group as a presented group, the map $f$, and the equivariant surjection $P^{ab}\to\Z^\T$;
\item for each exceptional type, finite certificates (commutant-scalar constituents and $\pm1$ witnesses);
\item the finite part of Lemma~\ref{lem:induced}: by a search-tree (orbit--stabiliser) certificate, every
reflection-preserving automorphism of $\rho(W)$ lies in $\rho(W)\langle\pi_0\rangle$, for all six types.
\end{itemize}
For each of the six types the final theorem has exactly four hypotheses, the cited inputs:
$W$ is finite; $P$ is characteristic \cite{CP03}; $P^{ab}\cong\Z^N$ as groups, with $N=|\T|$ (Brieskorn); and
the Franzsen--Howlett decomposition $\Aut(W)=R(W)\mathcal A(W)$ \cite[Thm.~31]{FH03}. The formalisation uses
computer-found constituents rather than the hand constructions of Section~\ref{sec:exceptional}. Type $D_5$ is
not formalised.

\section{The characteristic reflection lattice}\label{sec:lattice}

\begin{lemma}\label{lem:reid}
Let $\psi\in\Aut(G)$.
\begin{enumerate}[label=(\alph*)]
\item If $K\trianglelefteq G$ with $\psi(K)=K$, then $R(\psi)\ge R(\bar\psi)$.
\item If $G=\Z^m$ and $\psi$ is given by $B$, then $R(\psi)=|\coker(I-B)|$. This is $\infty$ if $B$ has eigenvalue $1$.
\item $R(\iota_g\circ\psi)=R(\psi)$, where $\iota_g(x)=gxg^{-1}$.
\item If $[G:H]<\infty$ and $\psi(H)=H$, then $R(\psi|_H)\le[G:H]\,R(\psi)$.
\end{enumerate}
\end{lemma}
\begin{proof}
(a) The projection maps classes onto classes.

(b) $y\sim x$ iff $y-x\in(I-B)\Z^m$.

(c) $x\mapsto xg$ maps $\iota_g\psi$-classes bijectively to $\psi$-classes.

(d) Write $G=\bigsqcup_iHg_i$. If $y=zx\psi(z)^{-1}\in H$ with $z=hg_i$, then $y$ is $\psi|_H$-conjugate to
$g_ix\psi(g_i)^{-1}$.
\end{proof}

\emph{Notation.} For a Coxeter graph with vertex set $S$:
\begin{itemize}
\item $A=\langle S\mid\Pi(s,t;m_{st})=\Pi(t,s;m_{st})\rangle$, where $\Pi(s,t;m)=sts\cdots$ has $m$ letters;
\item $W=A/\langle\!\langle s^2\rangle\!\rangle$, with quotient map $\mu\colon a\mapsto\bar a$, and $P=\ker\mu$;
\item $\T$ is the set of reflections, $\rho(w)(r)=wrw^{-1}$, and $\rho(w)$ also denotes the permutation matrix on
$\Z^\T$.
\end{itemize}
Since $\T$ generates $W$, $\ker\rho=Z(W)$. Each $\alpha\in\Aut(W)$ induces $\theta_\alpha$ on $\rho(W)$ by
$\theta_\alpha\rho(w)=\rho(\alpha w)$. For finite irreducible $W$ of rank $\ge3$, $Z(W)\subseteq\{\pm1\}$, and
$r\mapsto\rho(r)$ is injective on $\T$.

\begin{theorem}[{\cite[Thm.~1.4]{CP03}}]\label{thm:char}
If $A$ is irreducible of spherical type, then $P$ is characteristic in $A$.
\end{theorem}

This was extended to reducible types in \cite{FP06}. The proof in \cite{CP03} is computer-assisted in several
cases:
\begin{itemize}
\item for $H_4,F_4,E_6,E_7,E_8$ it uses GAP calculations \cite[Prop.~5.4, Cor.~5.5]{CP03};
\item for $H_3$ it uses one GAP calculation \cite[Lemma~5.1]{CP03};
\item among the classical types it uses GAP for $B_4,B_6,D_4,D_6$ \cite[Lemmas~3.5, 4.4]{CP03}.
\end{itemize}
We have not re-run these calculations. Each $\varphi\in\Aut(A)$ thus induces $\varphi_W\in\Aut(W)$ and
$\theta_\varphi:=\theta_{\varphi_W}$.

\begin{lemma}\label{lem:f}
For any Coxeter group, $s\mapsto(e_s,s)$ defines a homomorphism $f\colon A\to\Z^\T\rtimes W$, with
$(v,w)(v',w')=(v+w\cdot v',ww')$ and $w\cdot e_r=e_{wrw^{-1}}$.
\end{lemma}
\begin{proof}
For $\Pi(s,t;m)$, the $\Z^\T$-component is $\sum_{k=1}^me_{r_k}$ with $r_k=(st)^{k-1}s$. For $\Pi(t,s;m)$ it is
$\sum_ke_{r'_k}$ with $r'_k=(st)^{m-k}s$. Both sets equal $\{(st)^js:0\le j<m\}$. Compare \cite[Lemma~3.7]{CP03}.
\end{proof}

\begin{proposition}\label{prop:H1}
Let $A$ be of spherical type. Then $f(P)=2\Z^\T$, and $f$ induces an isomorphism $H_1(P;\Z)\cong2\Z^\T$ under which
conjugation by $a$ acts as the permutation $\rho(\bar a)$.
\end{proposition}
\begin{proof}
We have $f(as^2a^{-1})=(2e_{\bar as\bar a^{-1}},1)$, and these elements generate $f(P)$. Moreover $P\cong\pi_1$ of
the complexified reflection-arrangement complement \cite{Bri71}, whose $H_1$ is free on the meridians, one per
hyperplane \cite{Bri73,OT92}. So $H_1(P;\Z)$ has rank $|\T|$, and the surjection onto $2\Z^\T$ is an isomorphism.
\end{proof}

\begin{corollary}\label{cor:twist}
Let $M_\varphi$ be the matrix of $\varphi_*$ on $\Z^\T$. Then
$M_\varphi\rho(w)M_\varphi^{-1}=\theta_\varphi(\rho(w))$ and $M_{\iota_a\circ\varphi}=\rho(\bar a)M_\varphi$.
Hence $\theta_\varphi$ preserves the permutation character $\chi_\T(g)=\#\{r:g(r)=r\}$.
\end{corollary}

The image $f(A)\cong A/[P,P]$ is the crystallographic quotient of \cite{GGO17,Mar15}; see also \cite{DKNS24}.

\section{The fixed-vector criterion}\label{sec:criterion}

A $\rho(W)$-invariant subspace $U\subseteq\Q^\T$ is \emph{good} if $U\otimes\C$ is irreducible and occurs with
multiplicity one in $\C^\T$. A good $U$ is a full isotypic component.

\begin{lemma}\label{lem:schur}
If $U$ is good and $N\in GL(\Z^\T)$ commutes with $\rho(W)$, then $NU=U$ and $N|_U=\pm1$.
\end{lemma}
\begin{proof}
$NU\cong U$, so $NU=U$. By Schur's lemma $N|_U=c\in\Q^\times$. Applying $N^{\pm1}$ to a primitive vector of
$U\cap\Z^\T$ gives $c^{\pm1}\in\Z$.
\end{proof}

\begin{lemma}\label{lem:cert}
A nonzero invariant $U\subseteq\Q^\T$ is good in each of the following cases.
\begin{enumerate}[label=(\roman*)]
\item Every orbital matrix of $\rho(W)$ on $\T\times\T$ acts on $U$ by a scalar.
\item $\langle\chi_U,\chi_U\rangle=\langle\chi_U,\chi_\T\rangle=1$.
\item $U\otimes\C$ is irreducible and $\sum_r\dim(U\otimes\C)^{\mathrm{Stab}(r)}=1$, where $r$ runs over
representatives of the $\rho(W)$-orbits on $\T$ (Frobenius reciprocity).
\end{enumerate}
\end{lemma}
\begin{proof}
(i) The orbital matrices span the commutant $\mathcal C=\End_{\C W}(\C^\T)=\prod_iM_{m_i}(\C)$, where
$\C^\T=\bigoplus_iV_i\otimes\C^{m_i}$. By hypothesis $\mathcal C$ maps $U_\C=U\otimes\C$ into itself and acts on it
by scalars. The central idempotents $e_i\in\mathcal C$ act on $U_\C$ by $0$ or $1$ and sum to $1$, so exactly one
$e_i$ acts as $1$. Hence $U_\C\subseteq V_i\otimes\C^{m_i}$ meets only one isotypic component. Being
$W$-invariant, $U_\C=V_i\otimes L$ for some $0\ne L\subseteq\C^{m_i}$, and $M_{m_i}(\C)$ preserves $L$ and acts on it
by scalars. This forces $m_i=1$ and $L=\C$. (ii) and (iii) are standard.
\end{proof}

\begin{theorem}\label{thm:crit}
Let $A$ be irreducible of spherical type. Suppose that for every $\varphi\in\Aut(A)$ there is a permutation $\pi$
of $\T$ with $\pi\rho(w)\pi^{-1}=\theta_\varphi(\rho(w))$ for all $w$, such that one of the following holds.
\begin{enumerate}[label=(\roman*)]
\item\emph{(stable)} There is a good $U$ with $\pi U=U$ such that, for each $\varepsilon=\pm1$, some $h\pi$
($h\in\rho(W)$) has eigenvalue $\varepsilon$ on $U$.
\item\emph{(swap)} There are good $U_1\ne U_2$ with $\pi U_1=U_2$ such that, for each $\varepsilon=\pm1$, some
$(h\pi)^2$ has an eigenvector in $U_1$ with eigenvalue $\varepsilon$.
\end{enumerate}
Then $A$ has property $\Rinf$.
\end{theorem}

\begin{proof}
$N:=\pi^{-1}M_\varphi$ commutes with $\rho(W)$.

\emph{Case (i).} $N|_U=c=\pm1$. Choose $h=\rho(\bar a)$ and $0\ne v\in U$ with $h\pi v=cv$, and let
$\varphi'=\iota_a\circ\varphi$. Then $M_{\varphi'}v=h\pi Nv=v$.

\emph{Case (ii).} $N|_{U_i}=c_i$. Choose $h=\rho(\bar a)$ and $0\ne v\in U_1$ with $(h\pi)^2v=c_1c_2v$. With
$Q=h\pi$ we have $QU_1=U_2$, and $M_{\varphi'}^2v=QNQNv=v$. Since $M_{\varphi'}v\in U_2$ and $U_1\cap U_2=0$,
$u=v+M_{\varphi'}v$ is a nonzero fixed vector.

So $M_{\varphi'}$ fixes a nonzero vector of $\Z^\T$. By Lemma~\ref{lem:reid}(b),(a) and
Proposition~\ref{prop:H1}, $R(\varphi'|_P)=\infty$. Then $R(\varphi')=\infty$ by (d), and $R(\varphi)=R(\varphi')$
by (c).
\end{proof}

\begin{remark}[Relation to the crystallographic criteria]\label{rem:crystal}
The image $f(A)\cong A/[P,P]$ is an extension of $W$ by the lattice $2\Z^\T$. If $Z(W)=1$, it is a crystallographic
group with holonomy representation $\rho$. In that case Theorem~\ref{thm:crit} is essentially the
Dekimpe--Penninckx criterion \cite[Cor.~3.10]{DP11}, \cite[Thm.~3.3]{DKT19}: $R=\infty$ iff
$\det(1-AD)=0$ for some holonomy element $A$. It is applied to $f(A)$, with the multiplicity-one device of
\cite[Thm.~4.7]{DDP09} used to produce the eigenvector. If $Z(W)\ne1$, the holonomy action has kernel $Z(W)$, so
$f(A)$ is not crystallographic in the strict sense. Working directly with $P$ avoids this issue.
\end{remark}

\begin{remark}\label{rem:reduce}
\begin{itemize}
\item If $\pi$ works for $\theta$, then $g\pi$ works for $\iota_g\theta$, since $\{hg\pi\}=\{h\pi\}$. So one
$\theta$ per coset of $\Inn(\rho(W))$ suffices.
\item For inner $\theta$ take $\pi=1$. Then (i) holds for any good $U$ on which some reflection acts nontrivially.
This is automatic for a nontrivial good $U$ when all reflections are conjugate: the kernel of the action is normal,
so if it contained one reflection it would contain all of them.
\item If $\theta(\rho(\T))=\rho(\T)$, then $\pi(r)=\rho^{-1}\theta\rho(r)$ realises $\theta$.
\end{itemize}
\end{remark}

\section{The exceptional types}\label{sec:exceptional}

\subsection{Induced automorphisms}
Let $W$ be finite irreducible. Let $R(W)\le\Aut(W)$ be the reflection-preserving automorphisms and $\Gr(W)$ the
graph automorphisms. A \emph{central twist} $w\mapsto w\,c(w)$, with $c\colon W\to Z(W)$ a homomorphism, induces the
identity on $W/Z(W)$. Franzsen and Howlett prove the following (\cite[Thm.~31 and the discussion preceding it, pp.~329--330]{FH03}).
\begin{itemize}
\item For $W$ of type $B_n,D_n,E_6,E_7,E_8,F_4,H_3,H_4$, $\Aut(W)=R(W)\mathcal A(W)$, where $\mathcal A(W)$ consists
of central twists.
\item $R(W)=\Inn(W)\Gr(W)$ except for $H_3$ and $H_4$, where $[R(W):\Inn(W)]=2$.
\end{itemize}
In \cite{FH03} (in the published version as in the preprint) the proof of Theorem~31 is given as an outline, based on the sizes of the conjugacy classes of
involutions. (In item (9) of that theorem, for $F_4$, ``$\alpha_\chi$'' should read $\alpha_\zeta$.) For the six
exceptional types, Appendix~\ref{app:out} gives independent control of all of $\Aut(\rho(W))$: by exhaustive
computation for $E_6,E_7,F_4,H_3,H_4$, and by hand for $E_8$.

\begin{lemma}\label{lem:induced}
Let $\varphi\in\Aut(A)$, with $A$ of type $X$. Then $\theta_\varphi\in\Inn(\rho(W))\,\Theta_X$, where:
\begin{itemize}
\item $\Theta_X=\{1\}$ for $X=E_6,E_7,E_8,B_n$ ($n\ge3$);
\item $\Theta_{F_4}=\{1,\theta_L\}$, $\Theta_{H_3}=\{1,\theta_\Sigma\}$, $\Theta_{H_4}=\{1,\theta_\Psi\}$;
\item $\Theta_{D_n}=\{\theta_\gamma:\gamma\in\Gr(D_n)\}$.
\end{itemize}
Every element of $\Theta_X$ maps $\rho(\T)$ onto itself, so it is realised by a permutation of $\T$.
\end{lemma}

\begin{proof}
Central twists act trivially on $\rho(W)$, so $\theta_\varphi$ is induced by an element of $R(W)$.
\begin{itemize}
\item For $E_6$ (and $D_n$, $n$ odd) the graph automorphism is induced by conjugation by $w_0$, hence inner. For $E_7$, $E_8$ and
$B_n$, $\Gr$ is trivial.
\item For $F_4$, conjugation by the similarity $L(x)=(x_1+x_2,x_1-x_2,x_3+x_4,x_3-x_4)$ satisfies
$Ls_\alpha L^{-1}=s_{L\alpha}$. It is reflection-preserving and exchanges the two reflection classes, so it is not
inner. Since $[R(W):\Inn(W)]\le|\Gr(W)|=2$, it represents the nontrivial coset.
\item For $H_3$, resp.\ $H_4$, conjugation by $\Sigma$, resp.\ $\Psi$ (Propositions~\ref{prop:H3}
and~\ref{prop:H4}) is reflection-preserving and induces a non-inner automorphism of $\rho(W)$. It therefore
represents the nontrivial coset of $\Inn(W)$ in $R(W)$.
\end{itemize}
Finally, $\theta_r(\rho(\T))=\rho(\T)$ for $r\in R(W)$, by definition of $R(W)$.
\end{proof}

\subsection{Explicit good constituents}

\begin{proposition}[$E_6,E_7,E_8$]\label{prop:E}
Let $A_\perp$ be the adjacency matrix of the orthogonality graph on root lines.
\begin{enumerate}[label=(\alph*)]
\item $\rho(W)$ has exactly three orbitals on $\T\times\T$: equality, orthogonality and non-orthogonality.
\item The graph is strongly regular with parameters $(36,15,6,6)$, $(63,30,13,15)$, $(120,63,30,36)$. Its
nontrivial eigenvalues are $3^{15},(-3)^{20}$; $3^{35},(-5)^{27}$; $3^{84},(-9)^{35}$ (multiplicities as
exponents).
\item The two nontrivial eigenspaces are good, and every reflection acts nontrivially on them.
\end{enumerate}
\end{proposition}

\begin{proof}
By Steinberg's theorem \cite{Ste64}, the stabiliser of a root, or of a pair of roots, is the Weyl group of the
perpendicular root subsystem. The perpendicular subsystems are listed below. The root counts and ranks were
checked by machine (Appendix~\ref{app:comp}). Each pair (number of roots, rank) in the table occurs for exactly one
simply-laced root system, irreducible or not, of that rank, so the counts determine the types. These are the
classical facts that the centraliser of a reflection in $W(E_8)$, $W(E_7)$, $W(E_6)$ involves $W(E_7)$, $W(D_6)$,
$W(A_5)$.

\begin{center}\small
\begin{tabular}{lccc}
\toprule
 & $\alpha^\perp$ & $\{\alpha,\gamma\}^\perp$, $\alpha\perp\gamma$ & $\{\alpha,\beta\}^\perp$, $(\alpha,\beta)=1$\\
\midrule
$E_6$ & $A_5$ (30, rank 5) & $A_3$ (12, rank 3) & $A_2^2$ (12, rank 4)\\
$E_7$ & $D_6$ (60, rank 6) & $D_4A_1$ (26, rank 5) & $A_5$ (30, rank 5)\\
$E_8$ & $E_7$ (126, rank 7) & $D_6$ (60, rank 6) & $E_6$ (72, rank 6)\\
\bottomrule
\end{tabular}
\end{center}

$W$ is transitive on roots. Inside $W(\alpha^\perp)$, the orbits are as follows.
\begin{itemize}
\item The orbit of an orthogonal root has size $720/24=30$, $23040/384=60$, $2903040/23040=126$, i.e.\ all of
$\alpha^\perp$.
\item The orbit of a root $\beta$ with $(\alpha,\beta)=1$ has size $720/36=20$, $23040/720=32$,
$2903040/51840=56$, i.e.\ all such roots.
\end{itemize}
Passing to lines gives (a). Then $k$, $\lambda$, $\mu$ are half the root counts in the table, and (b) follows from
the standard formulas for strongly regular graphs.

By (a) the commutant is spanned by $I$, $A_\perp$ and $J$. These act by scalars on the nontrivial eigenspaces,
which are rational because the eigenvalues are integers. So Lemma~\ref{lem:cert}(i) applies, and
Remark~\ref{rem:reduce} gives the last assertion.
\end{proof}

\begin{proposition}[$H_3$]\label{prop:H3}
\begin{enumerate}[label=(\alph*)]
\item $\rho(W)\cong A_5$. The $15$ root lines form $5$ orthogonal triples (\emph{frames}), which correspond to the
Klein four-subgroups.
\item The frame module $U$ (functions constant on frames, with sum $0$) is a good $4$-dimensional constituent.
More precisely, writing irreducible $A_5$-modules by their dimensions,
\[\Q^\T\cong1\oplus U\oplus5\oplus5 .\]
\item For $\Sigma=\varsigma\circ\sigma$, where $\sigma$ is Galois conjugation $\sqrt5\mapsto-\sqrt5$ on coordinates
and $\varsigma$ swaps the last two coordinates, $\Sigma\Phi=\Phi$ and $\Sigma s_\alpha\Sigma^{-1}=s_{\Sigma\alpha}$. The
induced $\pi_\Sigma$ is not inner on $\rho(W)$.
\item For every outer $\theta$ there is $h$ such that $h\pi$ acts on the frames as a transposition. It then has
eigenvalues $+1$ (multiplicity 3) and $-1$ (multiplicity 1) on $U$.
\end{enumerate}
\end{proposition}

\begin{proof}
(a) Commuting reflections correspond to orthogonal lines, and the $15$ involutions of $A_5$ lie in $5$ Klein
four-subgroups, permuted as $A_5$ permutes $5$ points.

(b) $U$ is the augmentation module of a $2$-transitive action, hence absolutely irreducible. Since
$\Q^\T=\mathrm{Ind}_{V_4}^{A_5}1$, the multiplicity of an irreducible $\chi$ is $\dim\chi^{V_4}$. This equals
$1,0,0,1,2$ for $\chi$ of degree $1,3,3',4,5$, which gives the displayed decomposition.

(c) $\Sigma$ satisfies $(\Sigma x,\Sigma y)=\sigma(x,y)$. $\theta_\Sigma$ is not inner: on the rotation subgroup $W^+\cong\rho(W)$ it Galois-conjugates the character of
$\mathbb R^3$ on elements of order $5$, while inner automorphisms preserve characters. This is also checked by
machine.

(d) An outer automorphism of $A_5$ is conjugation by an odd permutation, and $\pi$ permutes the frames accordingly.
So $\{h\pi\}$ contains a transposition. In fact $\pi_\Sigma$ itself is a transposition on frames.
\end{proof}

\begin{proposition}[$F_4$]\label{prop:F4}
Each class of $12$ root lines (long, short) splits into $3$ orthogonal frames of $4$ lines. Let
$U_{\rm long},U_{\rm short}$ be the corresponding frame modules (sum zero). Then:
\begin{enumerate}[label=(\alph*)]
\item both are good and $2$-dimensional;
\item long reflections act on $U_{\rm long}$ as transpositions of frames and trivially on $U_{\rm short}$, and
vice versa;
\item $\pi_L$ exchanges $U_{\rm long}$ and $U_{\rm short}$, and $\pi_L^2=1$;
\item for a long simple reflection $s$, $(s\pi_L)^2=s\cdot\pi_Ls\pi_L$ has eigenvalue $-1$ on $U_{\rm short}$, and
$\pi_L^2$ gives eigenvalue $+1$.
\end{enumerate}
Hence case (ii) of Theorem~\ref{thm:crit} holds for $\theta_L$.
\end{proposition}

\begin{proof}
The long roots $\pm e_i\pm e_j$ form a $D_4$, whose frames correspond to the pairings $ij|kl$. The short lines
$e_i$ form one frame. The lines $\frac12(\pm1,\pm1,\pm1,\pm1)$ form two frames, according to the parity of the
number of minus signs. Each frame system is permuted through all of $S_3$, so each frame module is the standard
module of $S_3$.

\emph{Multiplicity one} (Lemma~\ref{lem:cert}(iii)):
\begin{itemize}
\item the stabiliser of the long line $e_1+e_2$ fixes one long frame and contains $s_{e_1-e_2}$, which swaps the
other two, so $\dim U_{\rm long}^{\rm Stab}=1$;
\item the stabiliser of the short line $e_1$ contains $\Sym\{2,3,4\}$, which induces $S_3$ on pairings, so
$\dim U_{\rm long}^{\rm Stab}=0$;
\item $U_{\rm short}$ follows by applying $\pi_L$.
\end{itemize}

(b) Check on $s_{e_1-e_2}$ and $s_{e_1}$, and use conjugacy.

(c) $L$ maps orthogonal frames to orthogonal frames and long to short lines, and $L^2=2I$.

(d) $\pi_Ls\pi_L$ is a short reflection.
\end{proof}

\begin{lemma}[Icosian model of $H_4$]\label{lem:icosian}
Halve the $120$ roots of $H_4$ (Appendix~\ref{app:comp}) and read the first coordinate as the real part. These
form a group $I$ of unit quaternions (the binary icosahedral group), and $K=I/\{\pm1\}\cong A_5$. Identify $\T$ with
$K$. Then:
\begin{enumerate}[label=(\alph*)]
\item the reflection in a unit root $r$ is $x\mapsto-r\bar xr$, and it acts on $K$ by $x\mapsto rx^{-1}r$;
\item $\rho(W)=\{x\mapsto ax^{\pm1}b:a,b\in K\}\cong(K\times K)\rtimes C_2$.
\end{enumerate}
\end{lemma}

\begin{proof}
(a) For $|r|=1$, $r\bar xr=2(x,r)r-x$, using $\bar xr+\bar rx=2(x,r)$.

(b) By (a), $\rho(W)$ is contained in the right-hand group, which has order $7200=|W|/2=|\rho(W)|$.

The closure of $I$ and the formulas are also checked by machine.
\end{proof}

\begin{proposition}[$H_4$]\label{prop:H4}
Let $V$ be the $4$-dimensional augmentation module of $K\cong A_5$ acting on its five Klein four-subgroups, and
let $U\subseteq\Q^K$ be the span of the matrix coefficients $x\mapsto\operatorname{tr}(X\lambda_V(x))$, so
$U\cong\End(V)$ has dimension $16$.
\begin{enumerate}[label=(\alph*)]
\item $U$ is good, and $\Q^\T\cong1+(3\otimes3)+(3'\otimes3')+16+25$.
\item Let $\Psi=\Omega\circ\sigma$, where $\sigma$ is Galois conjugation and $\Omega$ is the rotation $j\mapsto k$,
$k\mapsto-j$. Then $\Psi$ is a ring automorphism of $\HH_{\Q(\sqrt5)}$ with $\Psi(I)=I$. The induced permutation
$\pi_\Psi$ normalises $\rho(W)$, preserves reflections, and is not inner.
\item For every outer $\theta$ there is $h$ such that $h\pi$ acts on $U$ as conjugation by a transposition. This
has eigenvalues $+1$ (multiplicity $10$) and $-1$ (multiplicity $6$).
\end{enumerate}
\end{proposition}

\begin{proof}
(a) By Wedderburn, $\C[K]=\bigoplus_{V'}V'\otimes V'^*$ as a $K\times K$-module, each summand irreducible and
occurring once. The summand for $V$ is $U$, which is rational and stable under inversion. The decomposition follows
from the irreducible degrees $1,3,3,4,5$ of $A_5$.

(b) $\Psi$ commutes with quaternion conjugation, so $\Psi s_r\Psi^{-1}=s_{\Psi r}$. Closure $\Psi(I)=I$ is checked
exactly. If $\theta_\Psi$ were inner, then $\pi_\Psi\in\rho(W)$, because the centraliser of $\rho(W)$ in $\Sym(K)$
is trivial. Since $\pi_\Psi$ fixes $1$, it would be $x\mapsto ax^{\pm1}a^{-1}$. The sign $-$ is impossible, since $\pi_\Psi$ is an
automorphism of the non-abelian group $K$. The sign $+$ is impossible, since conjugation preserves real parts up to
sign, whereas $\Psi$ maps an element of order $5$ with real part $\tau/2$ to one with real part $-\tau^{-1}/2$
(here $\tau=\frac{1+\sqrt5}2$).

(c) We argue in three steps.
\begin{enumerate}[label=(\roman*)]
\item By Lemma~\ref{lem:induced} and Remark~\ref{rem:reduce}, for an outer $\theta$ we may take
$\pi\in\rho(W)\pi_\Psi$. So $\{h\pi:h\in\rho(W)\}=\{h\pi_\Psi\}$.
\item By (b), $\pi_\Psi=\psi$ is an outer automorphism of $K$.
\item Since $\Aut(A_5)\cong S_5$, acting on the five Klein four-subgroups, $\psi$ is conjugation by an odd
permutation $\delta'$. For $a\in K=A_5$, the element $h=(x\mapsto axa^{-1})$ lies in $\rho(W)$, and $h\psi$ is
conjugation by $a\delta'$. Choose $a$ with $a\delta'=\delta$ a transposition.
\end{enumerate}
The action of $K$ on $V$ extends to $S_5$, with $\lambda_V(\delta)$ having eigenvalues $1,1,1,-1$. Since
$h\pi_\Psi$ acts on $U\cong\End(V)$ as $X\mapsto\lambda_V(\delta)X\lambda_V(\delta)^{-1}$, the eigenvalue $+1$ has
multiplicity $9+1$ and $-1$ has multiplicity $6$.
\end{proof}

\begin{proof}[Proof of the Main Theorem]
By Lemma~\ref{lem:induced} and Remark~\ref{rem:reduce}, it suffices to treat $\Theta_X$.
\begin{itemize}
\item $E_n$: only the inner class occurs. Proposition~\ref{prop:E} supplies $U$, and Remark~\ref{rem:reduce}
applies, since all reflections are conjugate.
\item $H_3$, $H_4$: the inner class is covered by Remark~\ref{rem:reduce} (all reflections are conjugate and $U$ is
nontrivial). The outer class is covered by Propositions~\ref{prop:H3} and~\ref{prop:H4}: there $\pi U=U$, since $U$
is the only good constituent of its dimension, and case~(i) applies.
\item $F_4$: there are two classes of reflections, so the inner class is treated explicitly. Take $\pi=1$ and
$U=U_{\rm long}$; then $h=1$ gives $+1$, and a long reflection gives $-1$ by Proposition~\ref{prop:F4}(b), so case~(i)
holds. The outer class $\theta_L$ is covered by Proposition~\ref{prop:F4}, case~(ii).
\item $D_5$: Proposition~\ref{prop:BD}.
\end{itemize}
Apply Theorem~\ref{thm:crit}.
\end{proof}

\begin{table}[ht]\centering\small
\begin{tabular}{lrrlll}
\toprule
type & $|\T|$ & $|\rho(W)|$ & $\Theta_X$ & good constituent used (dim) & outer witness\\
\midrule
$E_6$ & 36 & 51840 & $1$ & $\mathbf{SRG}$ eigenspace (15; also 20) & ---\\
$E_7$ & 63 & 1451520 & $1$ & $\mathbf{SRG}$ eigenspace (35; also 27) & ---\\
$E_8$ & 120 & 348364800 & $1$ & $\mathbf{SRG}$ eigenspace (84; also 35) & ---\\
$F_4$ & 24 & 576 & $\{1,\theta_L\}$ & $\mathbf{U_{long}},\mathbf{U_{short}}$ (2+2) & swap: $\pi_L^2$, $(s\pi_L)^2$\\
$H_3$ & 15 & 60 & $\{1,\theta_\Sigma\}$ & $\mathbf{frame}$ module (4) & eigen-multiplicities 3, 1\\
$H_4$ & 60 & 7200 & $\{1,\theta_\Psi\}$ & $\mathbf{\End(V)}$ (16) & eigen-multiplicities 10, 6\\
\bottomrule
\end{tabular}
\caption{The data used for the Main Theorem. The dimensions agree with independent computer certificates
(Appendix~\ref{app:comp}).}
\end{table}

\section{Uniform proofs for the classical types}\label{sec:classical}

Except for $D_5$, these results are known \cite{CS22,SV24}. We include them to illustrate the method, and because
for $D_5$ the argument gives what appears to be the first proof (Remark~\ref{rem:D5}).

\begin{proposition}[$A_n$, $n\ge2$]\label{prop:An}
$A(A_n)$ has property $\Rinf$.
\end{proposition}
\begin{proof}
Here $\rho(W)=S_{n+1}$ acts on the transpositions, i.e.\ on the $2$-subsets. By Young's rule
$\Q^\T=S^{(n+1)}\oplus S^{(n,1)}\oplus S^{(n-1,2)}$ (for $n=2$ the last summand is absent). So $S^{(n,1)}$ is good.
For $n+1\ne6$, $\Aut(S_{n+1})=\Inn$. For $n+1=6$, an outer automorphism maps $(12)\mapsto(12)(34)(56)$, but
$\chi_\T$ takes the values $7$ and $3$ there, so it is not of the form $\theta_\varphi$ (Corollary~\ref{cor:twist}).
Apply Remark~\ref{rem:reduce}.
\end{proof}

For $B_n$ and $D_n$ let $u_{ij}=e_{[e_i+e_j]}-e_{[e_i-e_j]}$ ($i<j$) and $U^-=\mathrm{span}\{u_{ij}\}$. Let $E\le W$
be the group of sign changes: all of them for $B_n$, the even ones for $D_n$.

\begin{proposition}[$B_n$, $n\ge3$; $D_n$, $n\ge5$]\label{prop:BD}
$U^-$ is good, and $A(B_n)$ and $A(D_n)$ have property $\Rinf$.
\end{proposition}
\begin{proof}
\emph{Weights.} $E$ acts on $u_{ij}$ by $\chi_{ij}(\varepsilon)=\varepsilon_i\varepsilon_j$. These characters are
nontrivial and pairwise distinct on $E$. For $D_n$ this needs $n\ge5$: on even sign changes $\chi_{ij}=\chi_{kl}$
iff $\varepsilon_i\varepsilon_j\varepsilon_k\varepsilon_l$ is trivial or the product of all signs.

\emph{Goodness.} $E$ fixes $e_{[e_i+e_j]}+e_{[e_i-e_j]}$ and the short lines. So each $\chi_{ij}$ occurs once in
$\Q^\T$. Since $S_n\le W$ permutes the weights transitively, $U^-\otimes\C$ is irreducible.

\emph{Automorphisms.} By Lemma~\ref{lem:induced}, $\theta_\varphi$ is inner for $B_n$, and inner up to the diagram
involution for $D_n$ (for odd $n$ the latter is inner too). For the inner class, $\pi=1$ works: the reflection
$s_{e_n}$ ($B_n$), resp.\ $s_{e_{n-1}-e_n}$ ($D_n$), acts nontrivially on $U^-$. The diagram involution is realised
by negating $e_n$, which acts on $u_{ij}$ by $-1$ if $n\in\{i,j\}$ and by $+1$ otherwise. So case~(i) holds with
$h=1$.
\end{proof}

\begin{remark}[Type $D_5$]\label{rem:D5}
Proposition~\ref{prop:BD} with $n=5$ proves $\Rinf$ for $A(D_5)$. Its inputs are \cite{CP03}, \cite{FH03} (here
$\Aut(W(D_5))=R(W)\cong W$ by \cite[Thm.~31(3)]{FH03}) and the weight argument.

There is also a short derivation from Paris--Soroko's description $\Aut(A[D_5])=\Inn(A)\rtimes\langle\chi\rangle$,
with $\chi$ the inversion of the generators \cite[Thm.~1]{PS26}:
\begin{itemize}
\item $R(\iota_g)=R(\mathrm{id})=\infty$ by Lemma~\ref{lem:reid}(c), since $A^{ab}\cong\Z$ and Lemma~\ref{lem:reid}(a),(b)
give infinitely many conjugacy classes;
\item $R(\iota_g\circ\chi)=R(\chi)=\infty$ by Remark~\ref{rem:elem}(b).
\end{itemize}
This deduction is nearly immediate once \cite[Thm.~1]{PS26} is available, but we have not found it written down.
\cite{SV24} states $D_5$ as a conjecture, and \cite{PS26} does not mention $\Rinf$. We therefore record $D_5$ as part
of the Main Theorem.
\end{remark}

\begin{proposition}[$D_4$]\label{prop:D4}
In type $D_4$:
\begin{itemize}
\item $U^-$ is reducible: $J(u_{ij})=u_{kl}$ ($\{k,l\}$ complementary) commutes with $W(D_4)$, and $U^-$ is the sum
of two good $3$-dimensional constituents;
\item the frame module $U$ on the $3$ orthogonal frames $\{12|34,13|24,14|23\}$ is good of dimension $2$;
\item $A(D_4)$ has property $\Rinf$.
\end{itemize}
\end{proposition}
\begin{proof}
The $E$-fixed part of $\Q^\T$ is the permutation module on the $6$ pairs, namely $1+3+2$ for $S_4$. Its
$2$-dimensional summand is $U$, which does not occur in $U^-$. Every $\theta_\varphi$ preserves reflections, so
$\pi$ permutes the frames. $W$ induces all of $S_3$ on frames, so some $h\pi$ is a transposition there, with
eigenvalues $\pm1$ on $U$.
\end{proof}

\begin{remark}[Dihedral types]\label{rem:dihedral}
Let $m$ be odd. Then $Z(W)=1$, so $\rho$ is faithful and injective on $\T$, and $\Q^\T\cong\bigoplus_{d\mid m}U_d$ with $\dim U_d=|(\Z/d\Z)^\times|$, and $U_d$ is
absolutely irreducible iff $d\in\{1,3\}$.
\begin{itemize}
\item If $3\nmid m$ there is no nontrivial good constituent, and Theorem~\ref{thm:crit} cannot be applied. For $m$
prime, Schur's lemma gives only a unit of $\Q(\cos\frac{2\pi}m)$.
\item If $3\mid m$, then $U_3$ is the unique good constituent of dimension $2$, so $\pi U_3=U_3$. The group
$\rho(W)$ acts on $U_3$ as $S_3\le O(2)$. Since $\pi$ is orthogonal and normalises $S_3$, it acts as an element of
$S_3\times\{\pm1\}$. Hence some $h\pi$ acts as $\pm$(a reflection), with eigenvalues $\pm1$. So
Theorem~\ref{thm:crit} applies, using \cite{CP03} and that every automorphism of $W$ preserves reflections.
\end{itemize}
This was confirmed by computation for odd $m\le21$. We have not analysed even $m$. In any case $\Rinf$ for
$I_2(m)$ is known \cite{CS22}.
\end{remark}

\begin{remark}[Two elementary observations]\label{rem:elem}
\begin{enumerate}[label=(\alph*)]
\item If $\varphi$ fixes the generator $z\in\{\Delta,\Delta^2\}$ of $Z(A)$ \cite{BS72,Del72}, then $\varphi_*$ fixes
the nonzero image of $z$ in $H_1(A)$, so $R(\varphi)=\infty$. Only centre-inverting automorphisms need the
criterion.
\item For any Artin group that is not free abelian, the inversion $s\mapsto s^{-1}$ has $R=\infty$. Choose
$m_{st}\ge3$. Via $f$, $\max_r|v_r-v_{srs}|$ is an invariant of twisted classes with $W$-part $s$, and it
separates the elements $t^{2k}s$. Cf.\ \cite[Thm.~1]{Sor26}, \cite{Juh11}.
\end{enumerate}
\end{remark}

\appendix
\section{Computational cross-checks}\label{app:comp}

The following were checked by exact arithmetic (over $\Q$ and $\Q(\sqrt5)$). Programs, commands and expected
outputs are listed in \texttt{code/README.md} of the accompanying repository.

\emph{Coordinates.} For $E_8$ we use Bourbaki's simple roots in $\mathbb R^8$ with basis $e_1,\dots,e_8$:
\[\alpha_1=\tfrac12(e_1+e_8)-\tfrac12(e_2+\cdots+e_7),\qquad\alpha_2=e_1+e_2,\qquad\alpha_k=e_{k-1}-e_{k-2}\ \ (3\le k\le8).\]
$E_7$ and $E_6$ use $\alpha_1,\dots,\alpha_7$, resp.\ $\alpha_1,\dots,\alpha_6$, in $\mathbb R^8$. The other types:
\begin{itemize}
\item $F_4$: $e_2-e_3,\ e_3-e_4,\ e_4,\ \frac12(e_1-e_2-e_3-e_4)$;
\item $A_n,B_n,D_n$: standard simple roots;
\item $H_3$: roots $\pm2e_i$ and the even permutations of $(\pm\tau,\pm1,\pm\tau^{-1})$, with $\tau=\frac{1+\sqrt5}2$;
\item $H_4$: roots $\pm2e_i$, $(\pm1,\pm1,\pm1,\pm1)$ and the even permutations of $(\pm\tau,\pm1,\pm\tau^{-1},0)$.
\end{itemize}

\emph{Checks.}
\begin{enumerate}[label=(\arabic*)]
\item Proposition~\ref{prop:E}: SRG parameters, three orbitals, eigenspace dimensions, perpendicular root counts and
ranks.
\item Propositions~\ref{prop:H3}--\ref{prop:H4} and Lemma~\ref{lem:icosian}:
\begin{itemize}
\item frames, and character norms $\langle\chi,\chi\rangle=\langle\chi,\chi_\T\rangle=1$;
\item the images of $W$ on frames ($A_5$, $S_3$);
\item the action of $\pi_\Sigma$, $\pi_L$, $\pi_\Psi$, with eigenspace dimensions $(3,1)$, $(1,1)$, $(10,6)$;
\item the closure of $I$, the reflection formula, and $\rho(W)=\{x\mapsto ax^{\pm1}b\}$.
\end{itemize}
\item Propositions~\ref{prop:BD} and~\ref{prop:D4}:
\begin{itemize}
\item the weight argument for $B_3$--$B_6$ and $D_4$--$D_7$, including its failure for $D_4$;
\item the operator $J$ and the frame module of $D_4$;
\item witnesses for all six graph automorphisms of $D_4$.
\end{itemize}
\item Earlier, independent certificates for the six exceptional types:
\begin{itemize}
\item commutant certificates;
\item character certificates written as fresh code;
\item witnesses for abstractly computed outer automorphisms.
\end{itemize}
They give the same dimensions $15,20$; $27,35$; $35,84$; $4$; $2+2$; $16,25$.
\item The criterion (for automorphisms preserving $\chi_\T$) checked uniformly for $A_2$--$A_7$, $B_3$--$B_7$,
$D_4$--$D_7$, odd $I_2(m)$ with $m\le21$, and the exceptional types other than $E_8$.
\end{enumerate}

\section{Corroboration: all automorphisms of $W/Z(W)$}\label{app:out}

The argument above uses only automorphisms of $\rho(W)$ induced from $\Aut(W)$. As independent corroboration
of \cite[Thm.~31]{FH03} for the exceptional types (whose proof there is an outline), we record the stronger
statement about all of $\Aut(\rho(W))$:
\begin{itemize}
\item $\Out(\rho(W))=1$ for $E_6$, $E_7$, $E_8$;
\item $|\Out(\rho(W))|=2$ for $H_3$, $H_4$, $F_4$.
\end{itemize}
For $E_6,E_7,H_3,H_4,F_4$ this was verified by exhaustive computation in two ways: brute force over generator
images, and involution classes followed by orbits of Coxeter tuples. For $E_8$ the following argument by hand gives
it.

\begin{lemma}\label{lem:sign}
Let $\Phi$ be an irreducible simply-laced root system of rank $\ge3$ (so its diagram is a tree). Then every
$\theta\in\Aut(\rho(W))$ with $\theta\rho(\T)=\rho(\T)$ is inner up to a graph automorphism.
\end{lemma}
\begin{proof}
Write $\theta\rho(s_i)=\rho(s_{\beta_i})$; this uses that $\rho$ is injective on $\T$. Orders of products agree in $W$
and in $\rho(W)$. Indeed, for roots $\beta\ne\pm\gamma$ the element $s_\beta s_\gamma$ is a plane rotation fixing a
complement of dimension $\ge1$, so no nontrivial power of it lies in $Z(W)\subseteq\{\pm1\}$; in particular
$s_is_j\ne-1$. Hence $(\beta_i,\beta_j)=0$ or $\pm1$ according to the diagram. Choose the signs along the tree. Then $g\colon\alpha_i\mapsto\beta_i$ is an isometry, and
$g\Phi=(gWg^{-1})\{\beta_i\}\subseteq\Phi$. So $g\{\alpha_i\}$ is a simple system. As $W$ is simply transitive on simple systems, $g=w\gamma$ with $w\in W$ and
$\gamma$ an isometry permuting $\{\alpha_i\}$, i.e.\ $g\in W\rtimes\Gr=\Aut(\Phi)$ \cite{Hum90}. Compare
\cite[Thm.~18, Cor.~19]{FH03}.
\end{proof}

\begin{proposition}\label{prop:E8}
Every automorphism of $W(E_8)/\{\pm1\}$ is inner.
\end{proposition}
\begin{proof}
$G^+=W^+/\{\pm1\}$ is the commutator subgroup, of index $2$. An involution outside $G^+$ lifts to $w$ with $w^2=1$
and $\det w=-1$; if $w^2=-1$ then $\det w=1$. By \cite[Lemma~5]{Car72} (provable via Steinberg's theorem), $w$ is a product of $k=\dim\ker(w+1)$ reflections in
mutually orthogonal roots, and $k$ is odd. Replacing $w$ by $-w$, we may take $k\in\{1,3\}$. For $k=3$, the
$(-1)$-eigenspace of $w$ is spanned by the three roots, and contains only the $6$ roots $\pm\beta_i$: a root
$\sum c_i\beta_i$ has $2c_i=(\cdot,\beta_i)\in\{0,\pm1,\pm2\}$ and $\sum c_i^2=1$. So $w$ determines the triple,
and $-w$ (with $k=5$) is not of this form.

The products for $k=3$ form one class of size $240\cdot126\cdot60/48=37800$, by transitivity along
$E_8\supset E_7\supset D_6$. So the reflections form the unique class of size $120$ outside $G^+$. Hence every
automorphism preserves $\rho(\T)$, and Lemma~\ref{lem:sign} applies.
\end{proof}

\begin{thebibliography}{99}
\bibitem{AIMPL} AIM Problem List, \emph{Geometry and topology of Artin groups}, Problem 8.6, aimpl.org/geomartingp/8.
\bibitem{Bri71} E. Brieskorn, Die Fundamentalgruppe des Raumes der regul\"aren Orbits einer endlichen komplexen Spiegelungsgruppe, \emph{Invent.\ Math.} 12 (1971), 57--61.
\bibitem{Bri73} E. Brieskorn, Sur les groupes de tresses [d'apr\`es V.\,I. Arnol'd], S\'em.\ Bourbaki 1971/72, Exp.~401, LNM 317, Springer, 1973, 21--44.
\bibitem{BS72} E. Brieskorn, K. Saito, Artin-Gruppen und Coxeter-Gruppen, \emph{Invent.\ Math.} 17 (1972), 245--271.
\bibitem{CS22} M. Calvez, I. Soroko, Property $R_\infty$ for some spherical and affine Artin--Tits groups, \emph{J.\ Group Theory} 25 (2022), 1045--1054.
\bibitem{Car72} R.\,W. Carter, Conjugacy classes in the Weyl group, \emph{Compositio Math.} 25 (1972), 1--59.
\bibitem{CP03} A.\,M. Cohen, L. Paris, On a theorem of Artin, \emph{J.\ Group Theory} 6 (2003), 421--441; arXiv:math/0204242 (same theorem numbering).
\bibitem{DDP09} K. Dekimpe, B. De Rock, P. Penninckx, The $R_\infty$ property for infra-nilmanifolds, \emph{Topol.\ Methods Nonlinear Anal.} 34 (2009), 353--373.
\bibitem{DGO21} K. Dekimpe, D.\,L. Gon\c{c}alves, O. Ocampo, The $R_\infty$ property for pure Artin braid groups, \emph{Monatsh.\ Math.} 195 (2021), 15--33.
\bibitem{DKT19} K. Dekimpe, T. Kaiser, S. Tertooy, The Reidemeister spectra of low dimensional crystallographic groups, \emph{J.\ Algebra} 533 (2019), 353--375; arXiv:1709.07649.
\bibitem{DP11} K. Dekimpe, P. Penninckx, The finiteness of the Reidemeister number of morphisms between almost-crystallographic groups, \emph{J.\ Fixed Point Theory Appl.} 9 (2011), 257--283.
\bibitem{Del72} P. Deligne, Les immeubles des groupes de tresses g\'en\'eralis\'es, \emph{Invent.\ Math.} 17 (1972), 273--302.
\bibitem{DKNS24} N.\,K. Dhanwani, P. Kumar, T.\,K. Naik, M. Singh, Linear representations, crystallographic quotients, and twisted conjugacy of virtual Artin groups, arXiv:2409.10270.
\bibitem{FG10} A. Fel'shtyn, D.\,L. Gon\c{c}alves, Twisted conjugacy classes in symplectic groups, mapping class groups and braid groups, \emph{Geom.\ Dedicata} 146 (2010), 211--223.
\bibitem{FN16} A. Fel'shtyn, T. Nasybullov, The $R_\infty$ and $S_\infty$ properties for linear algebraic groups, \emph{J.\ Group Theory} 19 (2016), 901--921.
\bibitem{FT12} A. Fel'shtyn, E. Troitsky, Twisted conjugacy classes in residually finite groups, arXiv:1204.3175.
\bibitem{FT15} A. Fel'shtyn, E. Troitsky, Aspects of the property $R_\infty$, \emph{J.\ Group Theory} 18 (2015), 1021--1034.
\bibitem{FFW23} F. Fournier-Facio, R.\,D. Wade, Aut-invariant quasimorphisms on groups, arXiv:2211.00800.
\bibitem{FFIMSW25} F. Fournier-Facio, H. Iveson, A. Martino, W. Sgobbi, P. Wong, Property $R_\infty$ for groups with infinitely many ends, arXiv:2504.12002.
\bibitem{FP06} N. Franco, L. Paris, On a theorem of Artin, II, \emph{J.\ Group Theory} 9 (2006), 731--751.
\bibitem{FH03} W.\,N. Franzsen, R.\,B. Howlett, Automorphisms of nearly finite Coxeter groups, \emph{Adv.\ Geom.} 3 (2003), 301--338, \url{https://doi.org/10.1515/advg.2003.018}. Theorem numbering agrees with the authors' preprint, \url{https://www.maths.usyd.edu.au/u/bobh/fh02.pdf}.
\bibitem{GGO17} D.\,L. Gon\c{c}alves, J. Guaschi, O. Ocampo, A quotient of the Artin braid groups related to crystallographic groups, \emph{J.\ Algebra} 474 (2017), 393--423.
\bibitem{HL24} G. Hiss, R. Lutowski, The $R_\infty$-property of flat manifolds: toward the eigenvalue one property of finite groups, arXiv:2412.10206.
\bibitem{Juh11} A. Juh\'asz, Twisted conjugacy in certain Artin groups, in: \emph{Ischia Group Theory 2010}, World Sci., 2012, 175--195.
\bibitem{Kourovka} E.\,I. Khukhro, V.\,D. Mazurov (eds.), \emph{The Kourovka Notebook}, 21st ed., Novosibirsk, 2026 (Problem 19.28, as cited in \cite[\S6]{SV24}).
\bibitem{Mar15} I. Marin, Crystallographic groups and flat manifolds from complex reflection groups, \emph{Geom.\ Dedicata} 182 (2016), 233--247.
\bibitem{OT92} P. Orlik, H. Terao, \emph{Arrangements of Hyperplanes}, Springer, 1992.
\bibitem{PS26} L. Paris, I. Soroko, Automorphisms of the Artin group of type $D_5$, arXiv:2606.26082.
\bibitem{Sor26} I. Soroko, Property $R_\infty$ for Artin and Coxeter groups whose outer automorphisms are length-preserving, arXiv:2610.13904.
\bibitem{SV24} I. Soroko, N. Vaskou, Property $R_\infty$ for new classes of Artin groups, arXiv:2409.18123.
\bibitem{SV26} I. Soroko, N. Vaskou, Acylindrical hyperbolicity of automorphism groups and property $R_\infty$ for generalized Higman groups, arXiv:2604.27526.
\bibitem{Hum90} J.\,E. Humphreys, \emph{Reflection Groups and Coxeter Groups}, Cambridge Univ.\ Press, 1990.
\bibitem{Ste64} R. Steinberg, Differential equations invariant under finite reflection groups, \emph{Trans.\ Amer.\ Math.\ Soc.} 112 (1964), 392--400.
\end{thebibliography}

\end{document}
